% The Space CoBot is an approach to a microgravity UAV, capable of omnidirectional flying, to operate inside the International Space Station, as well as other microgravity atmosphered environments. In the development process, we selected the following guidelines: 

As starting point for the design of the aerial vehicle we considered
the following requirements:
\begin{enumerate}
\itemsep=0em
\item Safe operation;
\item Self-contained and simple propulsion system, with holonomic
  kinematics, that is, capable of motion in an arbitrary direction;
\item Autonomous Guidance, Navigation and Control (GN\&C);
\item Small dimensions;
\item Modularity;
\item Stackability and docking.
\end{enumerate}
An holonomic kinematics implies the decoupling between translational
and rotational motion. That is, the vehicle should not be required to
change attitude to move in a given direction. This requirement,
together with the one of small dimensions is particularly relevant
inside confined spaces such are those of space stations. We also
require the vehicle to have an autonomous GN\&C, thus simplifying the
deployment in already equipment-dense space environments. Modularity is
relevant to not only simplify maintenance, but also to allow the vehicle
to be extensible. Stackability and docking targets not only the dock
into a charging station, but also the possibility of parking them one
 on top of the others, thus saving space.

%%% Local Variables:
%%% mode: latex
%%% TeX-master: "main"
%%% End:
